\section{La experiencia de DJI: del control de vuelo a la confiabilidad de nivel de producto}

\subsection{La evolución del sistema de control de vuelo de DJI}

Yang Shuo (杨硕) empezó a participar en el desarrollo del sistema de control de vuelo de DJI desde 2013, durante sus estudios de posgrado. Muchos de los alumnos de su asesor, el profesor Li Zeqiang (李泽强), colaboraban de manera profunda con DJI, lo cual le dio a Yang Shuo una transición natural del ámbito académico al industrial.

Durante sus cinco años en DJI (2013--2018), Yang Shuo vivió de cerca la evolución completa del sistema de control de vuelo, desde lo básico hasta la madurez:
\begin{itemize}
    \item Participó y dirigió la reestructuración de los algoritmos del sistema de control de vuelo más temprano de DJI
    \item Trabajó en el rediseño de los algoritmos de control de bajo nivel y de la arquitectura del sistema de software
    \item Participó en varios ciclos de desarrollo de productos de consumo
    \item En la etapa final, estuvo a cargo del proyecto de la competencia de robótica RoboMaster
\end{itemize}

\subsection{La búsqueda de una confiabilidad comparable a la de Boeing}

Yang Shuo compartió un detalle que deja una impresión notable: en 2016, hubo una reunión interna en DJI en la que se anotaron en el pizarrón los competidores, y el nombre más grande de la lista era \textbf{Boeing}. En ese momento, la tasa de fallas del Phantom 4 era de aproximadamente uno en dos millones, mientras que Boeing lograba una tasa de uno en cinco millones. La meta de DJI era equipararse a la confiabilidad de los sistemas de nivel aeroespacial.

La lógica detrás de esta búsqueda es la siguiente: cuando se venden un millón o incluso dos millones de dispositivos voladores, aun con una tasa de fallas de uno en dos millones, la cantidad absoluta de accidentes sigue siendo considerable. Por eso es necesario seguir mejorando la confiabilidad de manera continua.

\begin{importantbox}{Lo que la confiabilidad de los drones enseña a la robótica}
Yang Shuo señala que la razón por la que DJI se propuso una meta de confiabilidad comparable a la de Boeing es que, tras el lanzamiento del Phantom 1 en 2013, se generaron ventas a gran escala, y al aumentar el número de usuarios comenzaron a presentarse accidentes, lo que obligó a la empresa a elevar la confiabilidad hasta un nivel aeroespacial. Esta experiencia tiene una relevancia directa para la industria de los robots humanoides: cuando en el futuro los robots de consumo alcancen la producción masiva, enfrentarán retos de confiabilidad similares. Las empresas de robótica actuales tienen dificultades incluso para lograr una tasa de fallas de uno en doscientos o uno en dos mil.
\end{importantbox}

\subsection{El foso competitivo de DJI: el poder de los detalles}

Yang Shuo considera que, 12 años después del lanzamiento del Phantom, el núcleo del foso tecnológico de DJI ya no es un avance técnico aislado, sino \textbf{la búsqueda extrema del detalle en todos los aspectos}. La manera de construir el control de vuelo de un dron, de ajustar la cámara, de diseñar la transmisión de video: todo esto ya cuenta con soluciones maduras dentro de la industria, pero al juntar todos los componentes, «el producto de nadie llega a ser tan bueno como el de DJI».

Esto ofrece una lección importante para la industria de la robótica: la convergencia tecnológica es inevitable —una tecnología madura siempre termina replicándose a gran escala—, pero la acumulación de detalles genera diferencias enormes. La complejidad de un robot supera en más de un orden de magnitud a la de un dron (cuatro motores frente a decenas de articulaciones), por lo que las diferencias en los detalles solo se amplifican aún más.

\subsection{La cadena completa del ciclo de vida del producto}

Yang Shuo enfatiza que el éxito de un producto de hardware va mucho más allá de la tecnología en sí misma, y requiere cubrir el ciclo de vida completo del producto:

\begin{enumerate}
    \item \textbf{Implementación técnica}: consolidar la tecnología básica (la mayoría de las empresas emergentes hoy se concentran en este paso)
    \item \textbf{Producción masiva}: pasar del prototipo a la fabricación a escala
    \item \textbf{Servicio posventa}: la atención de fallas una vez que el usuario tiene el producto en sus manos
    \item \textbf{Reparación}: el deterioro del desempeño tras un periodo de uso y el sistema de mantenimiento correspondiente
\end{enumerate}

Yang Shuo señala que la razón por la que Tesla es la empresa más cercana al nivel de Boeing en confiabilidad de robots es precisamente que su negocio de robótica heredó la infraestructura completa ya establecida por su línea de negocio automotriz: desde el despliegue, el hardware y la recolección de datos hasta las pruebas, todo cuenta con sistemas y procesos muy bien desarrollados.

\subsection{Resumen de la sección}

La experiencia en DJI le dio a Yang Shuo una perspectiva completa que va de lo académico a lo industrial. La búsqueda de una confiabilidad comparable a la de Boeing, la atención extrema al detalle y una comprensión profunda del ciclo de vida completo del producto conforman la metodología central de su emprendimiento posterior. Una idea central es esta: el foso tecnológico no reside en un avance aislado, sino en la optimización continua y la acumulación sistemática de innumerables detalles.
